% !PRINTED 354
\weberpage{354}

Wir wollen dies auf den Fall b) anwenden; dabei ist zu unterscheiden, ob $\lambda$ kleiner oder grösser als $2^m$ ist, weil davon abhängt, ob $\Theta$ im ersten oder zweiten Octanten liegt. Es sind also wieder zwei Fälle zu unterscheiden:
\[
\alpha)\quad \lambda < 2^m,
\]
\[
\lambda = \frac{b^m}{a^n} = \frac{c^m}{a^{m+n}} = \frac{b^{m+n}}{c^n},
\]
\[
(5)\quad x^{m+n} = gc, \quad x^n = eb, \quad x^m = \frac{ga}{e}.
\]
\[
\beta)\quad \lambda > 2^m,
\]
\[
\lambda = a^{m+n}b^m = a^n c^m = \frac{c^{m+n}}{b^n},
\]
\[
x^{m+n} = gb, \quad x^n = ec, \quad x^m = \frac{g}{ea}.
\]

Im Falle $\alpha)$ würde man also
\[
(6)\quad \log\lambda = mB - nA
\]
setzen und danach aus der Tafel die zusammengehörigen Werthe von $A$ und $B$ aufsuchen. Hat man diese gefunden, so ergiebt sich
\[
(7)\quad m\log x = A + \log g - \log e.
\]
Um den Gebrauch dieser Formeln an einem Beispiel zu erläutern, wollen wir die Gleichung betrachten:
\[
x^3 - 2x - 2 = 0,
\]
also $e = g = 2$, $m = 1$, $n = 2$, $\lambda = \frac{1}{2}$ setzen. Die Formel (6) giebt also
\[
(8)\quad 2A - B = 0{,}3010300
\]
und (7)
\[
(9)\quad \log x = A.
\]

Um den ersten Näherungswerth von $A$ zu finden, benutzt man einen kleinen Auszug aus der Tafel:
\[
\left|
\begin{array}{ll}
A = 0, & B = 0{,}301 \\
A = 0{,}1, & B = 0{,}254 \\
A = 0{,}2, & B = 0{,}212 \\
A = 0{,}3, & B = 0{,}176
\end{array}
\right|.
\]
Es ergiebt sich daraus:
\[
\begin{array}{c|c|c}
A & 2A - B & \text{Fehler} \\
\hline
0{,}2 & 0{,}188 & -0{,}113 \\
0{,}3 & 0{,}424 & +0{,}123
\end{array}
\]
